%% ****** Start of file apstemplate.tex ****** %
%%
%%
%%   This file is part of the APS files in the REVTeX 4.2 distribution.
%%   Version 4.2a of REVTeX, January, 2015
%%
%%
%%   Copyright (c) 2015 The American Physical Society.
%%
%%   See the REVTeX 4 README file for restrictions and more information.
%%
%
% This is a template for producing manuscripts for use with REVTEX 4.2
% Copy this file to another name and then work on that file.
% That way, you always have this original template file to use.
%
% Group addresses by affiliation; use superscriptaddress for long
% author lists, or if there are many overlapping affiliations.
% For Phys. Rev. appearance, change preprint to twocolumn.
% Choose pra, prb, prc, prd, pre, prl, prstab, prstper, or rmp for journal
%  Add 'draft' option to mark overfull boxes with black boxes
%  Add 'showkeys' option to make keywords appear
%\documentclass[aps,prl,preprint,groupedaddress]{revtex4-2}
\documentclass[aps,prl,twocolumn,groupedaddress]{revtex4-2}
%\documentclass[aps,prl,preprint,superscriptaddress]{revtex4-2}
%\documentclass[aps,prl,reprint,groupedaddress]{revtex4-2}

% You should use BibTeX and apsrev.bst for references
% Choosing a journal automatically selects the correct APS
% BibTeX style file (bst file), so only uncomment the line
% below if necessary.
%\bibliographystyle{apsrev4-2}

\usepackage{graphicx}
\usepackage{epstopdf}
\usepackage{caption}
%\usepackage{amsmath}% http://ctan.org/pkg/amsmath
%\usepackage{amsthm}
%\usepackage{amsfonts}
%\usepackage{subfigure}
%\usepackage{hhline}
%\usepackage[miktex]{gnuplottex}
%\usepackage{xcolor}
\usepackage{amssymb}
\usepackage{amsmath}
\usepackage{color}
\usepackage{hyperref}
%\usepackage[percent]{overpic}
\usepackage{tikz}
\usepackage{mathrsfs}
\usepackage{wasysym}
\usepackage{tikz-cd}
%\usepackage{linearb}
%\usepackage{wsuipa}
\usepackage{calc}
\usepackage{tipa}

%\usepackage{stix} %\fisheye
\usepackage{stackengine,scalerel}

% so sections, subsections, etc. become numerated.
\setcounter{secnumdepth}{3}

\newcommand{\FF}{\mathbb{F}}
\newcommand{\NN}{\mathbb{N}}
\newcommand{\ZZ}{\mathbb{Z}}
\newcommand{\QQ}{\mathbb{Q}}
\newcommand{\RR}{\mathbb{R}}
\newcommand{\CC}{\mathbb{C}}
\newcommand{\II}{\mathbb{I}}
\newcommand{\GG}{\mathbb{G}}

\DeclareMathOperator*{\argmax}{arg\,max}
\DeclareMathOperator*{\argmin}{arg\,min}
\newcommand{\avrg}[1]{\left\langle #1 \right\rangle}
\newcommand{\nelta}{\bar{\delta}}
\newcommand{\bra}[1]{\left\langle #1\right|}
\newcommand{\ket}[1]{\left| #1 \right\rangle}
\newcommand{\sbra}[1]{\langle #1|}
\newcommand{\sket}[1]{| #1 \rangle}
\newcommand{\bek}[3]{\left\langle #1 \right| #2 \left| #3 \right\rangle}
\newcommand{\sbek}[3]{\langle #1 | #2 | #3 \rangle}
\newcommand{\braket}[2]{\left\langle #1 \middle| #2 \right\rangle}
\newcommand{\ketbra}[2]{\left| #1 \middle\rangle \middle\langle #2  \right|}
\newcommand{\sbraket}[2]{\langle #1 | #2 \rangle}
\newcommand{\sketbra}[2]{| #1 \rangle  \langle #2 |}
\newcommand{\norm}[1]{\left\lVert#1\right\rVert}
\newcommand{\snorm}[1]{\lVert#1\rVert}
\newcommand{\bvec}[1]{\boldsymbol{\mathsf{#1}}}
\newcommand{\bcov}[1]{\boldsymbol{#1}}
\newcommand{\bdua}[1]{\boldsymbol{\check{#1}}}
%\newcommand{\bdov}[1]{\boldsymbol{\breve{#1}}}
\newcommand{\bdov}[1]{\breve{#1}}
%\newcommand{\bten}[1]{\boldsymbol{\mathfrak{#1}}}
\newcommand{\bten}[1]{\boldsymbol{\mathfrak{#1}}}
\newcommand{\forany}{\tilde{\forall}}
\newcommand{\qed}{$\overset{\circ}{.}\;$}
\newcommand{\bigo}{\mathcal{O}}
\newcommand{\spn}{\mathrm{spn}\,}
\newcommand{\krn}{\mathrm{krn}\,}
\newcommand{\biy}{\leftrightarrow}
\newcommand{\ugh}{\mathbin{\textlinb{\BPwheel}}}

\newcommand\bigeye{\ensurestackMath{\stackinset{c}{}{c}{-.3pt}%
  {\bullet}{\scriptstyle\bigcirc}}}
\newcommand\eye{\scalerel*{\bigeye}{x}}
%\newcommand*{\fisheye}{%
%    \mathbin{%
%        \ooalign{$\circledcirc$\cr\hidewidth$\bullet$\hidewidth}%
%    }%
%}

% double angle-bracket
% https://tex.stackexchange.com/questions/79657/how-to-get-double-angle-bracket-without-using-mnsymbol-package
\makeatletter
\newsavebox{\@brx}
\newcommand{\llangle}[1][]{\savebox{\@brx}{\(\m@th{#1\langle}\)}%
  \mathopen{\copy\@brx\mkern2mu\kern-0.9\wd\@brx\usebox{\@brx}}}
\newcommand{\rrangle}[1][]{\savebox{\@brx}{\(\m@th{#1\rangle}\)}%
  \mathclose{\copy\@brx\mkern2mu\kern-0.9\wd\@brx\usebox{\@brx}}}
\makeatother

% https://tex.stackexchange.com/questions/297362/bar-through-letter-in-mathmode
%\usepackage{calc}
\newsavebox\CBox
\newcommand\hcancel[2][0.5pt]{%
  \ifmmode\sbox\CBox{$#2$}\else\sbox\CBox{#2}\fi%
  \makebox[0pt][l]{\usebox\CBox}%  
  \rule[0.77\ht\CBox-#1/2]{\wd\CBox}{#1}}
\newcommand{\cb}{\hcancel{b}}
\newcommand{\cd}{\hcancel{d}}
\newcommand{\ct}{\hcancel{\partial}}
%\newcommand{\supp}{\mathsf{supp}}
\newcommand{\supp}{\mathrm{supp}}
\newcommand{\csupp}{\overline{\mathrm{supp}}}
\newcommand{\supr}[2]{\underset{#1}{\text{supr}}}

\begin{document}

% Use the \preprint command to place your local institutional report
% number in the upper righthand corner of the title page in preprint mode.
% Multiple \preprint commands are allowed.
% Use the 'preprintnumbers' class option to override journal defaults
% to display numbers if necessary
%\preprint{}

%Title of paper
\title{
M\'etodos Matem\'aticos de la F\'isica II\\
Part II: Analysis of Banach and Hilbert spaces
}

% repeat the \author .. \affiliation  etc. as needed
% \email, \thanks, \homepage, \altaffiliation all apply to the current
% author. Explanatory text should go in the []'s, actual e-mail
% address or url should go in the {}'s for \email and \homepage.
% Please use the appropriate macro foreach each type of information

% \affiliation command applies to all authors since the last
% \affiliation command. The \affiliation command should follow the
% other information
% \affiliation can be followed by \email, \homepage, \thanks as well.
\author{Juan I. Perotti}
\email[]{juan.perotti@unc.edu.ar}
%\homepage[]{Your web page}
%\thanks{}
%\altaffiliation{}
%\affiliation{}
\affiliation{Instituto de F\'isica Enrique Gaviola (IFEG-CONICET), Ciudad Universitaria, 5000 C\'ordoba, Argentina}
\affiliation{Facultad de Matem\'atica, Astronom\'ia, F\'isica y Computaci\'on, Universidad Nacional de C\'ordoba, Ciudad Universitaria, 5000 C\'ordoba, Argentina}

% \author{Colaborador}
% \affiliation{Instituto de F\'isica Enrique Gaviola (IFEG-CONICET), Ciudad Universitaria, 5000 C\'ordoba, Argentina}
% \affiliation{Facultad de Matem\'atica, Astronom\'ia, F\'isica y Computaci\'on, Universidad Nacional de C\'ordoba, Ciudad Universitaria, 5000 C\'ordoba, Argentina}
% \email[E-mail: ]{xxx.yyy@unc.edu.ar }

%Collaboration name if desired (requires use of superscriptaddress
%option in \documentclass). \noaffiliation is required (may also be
%used with the \author command).
%\collaboration can be followed by \email, \homepage, \thanks as well.
%\collaboration{Claudio J. Tessone}
%\noaffiliation

\date{\today}

\begin{abstract}
Bla bla bla... resumen...
\end{abstract}

% insert suggested keywords - APS authors don't need to do this
%\keywords{}

%\maketitle must follow title, authors, abstract, and keywords
\maketitle
\onecolumngrid

% https://chatgpt.com/share/6a0cb76f-d84c-83e9-abd4-078d5e57a990

\section{WHY BANACH AND HILBERT SPACES MATTER?}

In Part I, metric spaces, normed vector spaces, Banach spaces, inner product spaces, and Hilbert spaces were introduced. However, the true importance of Banach and Hilbert spaces was intentionally postponed.
\begin{itemize}
\item Why are normed vector spaces not enough?
\item Why are inner product spaces not enough?
\end{itemize}
The answer is essentially related to the convergence of infinite limiting processes.

In many areas of mathematics and physics, one works with:
\begin{itemize}
\item infinite sequences,
\item infinite series,
\item approximation procedures,
\item Fourier expansions,
\item variational methods, etc.
\end{itemize}
In such contexts, it is important that the limits produced by these processes remain inside the space itself.
Banach spaces guarantee that limits associated with norm convergence remain within the space. 

Hilbert spaces additionally provide a geometric structure based on orthogonality, allowing generalized Fourier expansions and orthogonal projections.
This part, Part II, develops these ideas.

\section{TOPOLOGICAL STRUCTURE OF METRIC SPACES}

\subsection{Open Balls}

Let $(X,d)$ be a metric space and let $x\in X$.
The open ball centered at $x$ with radius $r>0$ is defined by
$$
B_r(x):={y\in X:d(x,y)<r}.
$$
Open balls describe the local structure induced by the metric.

\subsection{Open Sets}

A subset $U\subseteq X$ is called open if for every $x\in U$, there exists $r>0$ such that
$$
B_r(x)\subseteq U.
$$
Intuitively, every point of an open set can be slightly perturbed while remaining inside the set.

\subsection{Closed Sets}

A subset $F\subseteq X$ is called closed if its complement
$$
X\setminus F
$$
is open.
Closed sets are important because they contain the limits of convergent sequences.

\subsection{Limit Points}

Let $A\subseteq X$.
A point $x\in X$ is called a limit point of $A$ if every open ball centered at $x$ contains points of $A$ different from $x$ itself.
That is,
$$
\forall r>0,\quad
(B_r(x)\setminus{x})\cap A\neq\emptyset.
$$

\subsection{Closure}
The closure of a subset $A\subseteq X$, denoted by $\overline{A}$, is the set obtained by adding to $A$ all of its limit points.
Equivalently, $\overline{A}$ is the smallest closed subset containing $A$.

\subsection{Dense Subsets}

A subset $A\subseteq X$ is called dense in $X$ if
$$
\overline{A}=X.
$$
This means that every point of $X$ can be approximated arbitrarily well by points of $A$.
Density is fundamental in approximation theory and Fourier analysis.



\subsection{Separable Spaces}

A metric space $X$ is called separable if it contains a countable dense subset.
That is, there exists a countable subset
$$
D=\{x_1,x_2,\dots\}\subseteq X
$$
such that
$$
\overline{D}=X.
$$
Intuitively, separability means that every element of the space can be approximated arbitrarily well using only countably many elements.
This property is fundamental in approximation theory, Fourier analysis, and quantum mechanics.
Examples of separable spaces include:
\begin{enumerate}
\item The $n$-th dimensional euclidean space $\mathbb{R}^n$.
Basically, it is separable because $\QQ$ is countable and dense in $\RR$.
\item The space $C[a,b]$ of continuous functions with domian $[a,b]\subset \RR$.
See appendix~\ref{app:B}.
\item The space $L^2[a,b]$ of square integrable functions with domain $[a,b]\subset \RR$.
\end{enumerate}

\subsection{Compact Sets}

A subset $K\subseteq X$ is called compact if every open cover of $K$ admits a finite subcover.
In metric spaces, compactness is equivalent to sequential compactness.

\subsection{Sequential Compactness}

A subset $K\subseteq X$ is sequentially compact if every sequence in $K$ possesses a convergent subsequence whose limit belongs to $K$.
Compactness plays a central role in existence theorems and approximation methods.

\subsection{Intuition behind Compactness}

To develop some intuition, let us state the {\em Heine-Borel theorem}.
It says that, for any $K\subseteq \RR^n$,
$$
K \text{ is compact}
\,\,\,
\Longleftrightarrow
\,\,\,
K \text{ is closed and bounded.}
$$
Thus:
\begin{itemize}
\item bounded = cannot escape to infinity,
\item closed = contains its limit points.
\end{itemize}
Together:
\begin{itemize}
\item sequences cannot diverge,
\item limits cannot leave the set.
\end{itemize}
So compactness means:
\begin{quotation}
every infinite process remains controlled.
\end{quotation}

\section{INFINITE SERIES AND BANACH SPACES}

\subsection{Infinite Series in Normed Spaces}

Let $(V,|\cdot|)$ be a normed vector space and let $\{v_n:n\in\mathbb{N}\}$ be a sequence in $V$.
An infinite series is an expression of the form
$$
\sum_{n=1}^\infty v_n.
$$
The series is said to converge if the sequence of partial sums
$$
S_N:=\sum_{n=1}^N v_n
$$
converges in norm.
That is, there exists $S\in V$ such that
$$
\lim_{N\to\infty}|S_N-S|=0.
$$

\subsection{The Cauchy Criterion for Series}

A series converges if and only if the sequence of partial sums is a Cauchy sequence.
That is,
$$
\forall\varepsilon>0,\quad
\exists N\in\mathbb{N}
\text{ such that }
m,n\geq N
\Rightarrow
|S_n-S_m|<\varepsilon.
$$
However, a Cauchy sequence does not necessarily converge in an arbitrary normed vector space.
This is precisely why completeness becomes important.

\subsection{Why Normed Spaces Are Not Enough}

In an incomplete normed vector space, a convergent approximation process may produce a limit that lies outside the space itself.
Thus, one may construct:
\begin{itemize}
\item convergent sequences,
\item convergent series,
\item approximation procedures,
\end{itemize}
whose limits do not belong to the original space.
This is mathematically inconvenient and physically undesirable.

\subsection{Banach Spaces}

Banach spaces solve this problem.
Because Banach spaces are complete, every Cauchy sequence converges to an element of the space itself.
Consequently:
\begin{itemize}
\item convergent infinite series remain inside the space,
\item approximation procedures remain meaningful,
\item limiting processes are internally consistent.
\end{itemize}
This is one of the main reasons Banach spaces are fundamental in analysis.

\subsection{Closed Subspaces of Banach Spaces}

Let $V$ be a Banach space and let $W\subseteq V$ be a closed vector subspace.
Then $W$ is itself a Banach space under the induced norm.
The closedness condition is essential because limits of convergent sequences in $W$ must remain inside $W$.

\section{ORTHOGONALITY AND HILBERT SPACES}

\subsection{Orthogonality}

Let $(H,\langle\cdot,\cdot\rangle)$ be an inner product space.
Two vectors $u,v\in H$ are said to be orthogonal if
$$
\langle u,v\rangle=0.
$$
Orthogonality generalizes the notion of perpendicularity from Euclidean geometry.

\subsection{Orthonormal Sets}

A subset $\{e_i:i\in I\}\subseteq H$ is called orthonormal if
$$
\langle e_i,e_j\rangle=
\delta_{ij}.
$$
That is:
\begin{itemize}
\item each vector has norm one,
\item different vectors are orthogonal.
\end{itemize}

\subsection{Fourier Coefficients}

Let $\{e_n\}_{n=1}^\infty$ be an orthonormal set and let $v\in H$.
The numbers
$$
c_n:=\langle e_n,v\rangle
$$
are called the Fourier coefficients of $v$ relative to the orthonormal set.
They measure the components of $v$ along the orthogonal directions $e_n$.

\subsection{Fourier Expansions}

One may attempt to reconstruct $v$ through the infinite series
$$
\sum_{n=1}^\infty
\langle e_n,v\rangle e_n.
$$
This is a generalized Fourier expansion.
However, such an expansion only becomes fully meaningful if the series converges within the space itself.
This is why completeness is essential.

\subsection{Bessel Inequality}

Let $\{e_n\}_{n=1}^\infty$  be an orthonormal set in a Hilbert space $H$.
Then, for every $v\in H$,
$$
\sum_{n=1}^\infty
|\langle e_n,v\rangle|^2
\leq
|v|^2.
$$
This inequality guarantees that the Fourier coefficients are square summable.

\subsection{Parseval Identity}

If the orthonormal set forms a complete orthonormal basis, then
$$
|v|^2
=
\sum_{n=1}^\infty
|\langle e_n,v\rangle|^2.
$$
This identity generalizes the Pythagorean theorem to infinite-dimensional spaces.

\subsection{Why Inner Product Spaces Are Not Enough?}

An arbitrary inner product space may fail to contain the limits of orthogonal expansions.
Thus, generalized Fourier series may converge to objects lying outside the space.
Hilbert spaces solve this problem because they are complete.
Therefore:
\begin{itemize}
\item orthogonal expansions converge inside the space,
\item generalized Fourier analysis becomes internally consistent,
\item infinite-dimensional Euclidean geometry becomes possible.
\end{itemize}
This is one of the main reasons Hilbert spaces are fundamental in mathematics and physics.

\section{ORTHOGONAL PROJECTIONS AND BEST APPROXIMATION}

\subsection{Orthogonal Complement}

Let $M\subseteq H$ be a subset of a Hilbert space.
The orthogonal complement of $M$ is defined by
$$
M^\perp
:=
{
v\in H:
\langle v,m\rangle=0
\text{ for all }m\in M
}.
$$

\subsection{Projection Theorem}

Let $M$ be a closed subspace of a Hilbert space $H$.
Then every vector $v\in H$ can be uniquely decomposed as
$$
v=m+w,
$$
where
\begin{itemize}
\item $m\in M$,
\item $w\in M^\perp$.
\end{itemize}
Equivalently,
$$
H=M\oplus M^\perp.
$$

\subsection{Best Approximation Theorem}

Let $M$ be a closed subspace of a Hilbert space $H$.
For every $v\in H$, there exists a unique vector $m\in M$ such that
$$
|v-m|
=
\min_{u\in M}|v-u|.
$$
Thus, orthogonal projections provide optimal approximations.

\subsection{Importance of Completeness}

The projection theorem fundamentally depends on completeness.
Without completeness:
\begin{itemize}
\item minimizing sequences may fail to converge,
\item orthogonal projections may fail to exist,
\item approximation procedures may break down.
\end{itemize}
Hilbert spaces guarantee that these geometric constructions remain meaningful.

\section{LINEAR OPERATORS}

\subsection{Linear Operators}

Let $V$ and $W$ be vector spaces.
A function
$$
T:V\to W
$$
is called linear if
$$
T(\alpha u+\beta v)
=
\alpha T(u)+\beta T(v)
$$
for all $u,v\in V$ and all scalars $\alpha,\beta$.

\subsection{Bounded Operators}

Let $V$ and $W$ be normed vector spaces.
A linear operator $T:V\to W$ is called bounded if there exists $C>0$ such that
$$
|T(v)|
\leq
C|v|
$$
for every $v\in V$.

\subsection{Continuity of Linear Operators}

A linear operator between normed vector spaces is continuous if and only if it is bounded.
Thus, boundedness is the natural notion of regularity for linear operators in functional analysis.

\subsection{Operator Norm}

The operator norm of a bounded linear operator $T$ is defined by
$$
|T|
:=
\sup_{|v|\leq1}
|T(v)|.
$$
The operator norm measures the maximal amplification produced by the operator.

\subsection{Final Remarks}

Banach and Hilbert spaces provide the natural setting for infinite-dimensional analysis.
Banach spaces guarantee the stability of limiting processes under norm convergence.
Hilbert spaces additionally provide a geometric structure based on orthogonality, projections, and generalized Fourier expansions.
These structures form the mathematical foundation of large areas of modern mathematics and theoretical physics.

\appendix

\section{
$C[a,b]$ is a Banach space
\label{app:A}
}

Let us show that $C[a,b]$ is a Banach space under the {\em uniform norm}.
$$
|f|_{\infty}
:=
\supr{x\in [a,b]}\,\,|f(x)|.
$$
For this, we should show that $C[a,b]$ is complete under such norm.
In other words, we should show that every Cauchy sequence in $C[a,b]$ converges to a function in $C[a,b]$.

That $\{f_n:n\in N\}$ is a Cauchy sequence under the uniform norm means that, for any $\epsilon>0$, there exists an $N\in \NN$ such that
$$
|f_n-f_m|_{\infty}<\epsilon
$$
for every $n,m\geq N$.
This is equivalent to say that
$$
|f_n(x)-f_m(x)|<\epsilon
$$
holds for every $x\in [a,b]$ and every $n,m\geq N$.
As a consequence, for any $x\in [a,b]$, the sequence $\{f_n(x):n\in \NN\}$ is Cauchy and, therefore, it is also convergent because $\RR$ is complete.
Let $f(x)$ denote the value to which $\{f_n(x):n\in \NN\}$ converges.
Now, it remains to show that:
\begin{enumerate}
\item the sequence $\{f_n(x):n\in \NN\}$ converges to $f$, and
\item that the function $f:[a,b]\ni x\to f(x)\in \RR$ belongs to $C[a,b]$ or, in other words, that $f$ is continuous in $[a,b]$.
\end{enumerate}

The condition 1. follows easily because the fact that $|f_n(x)-f_m(x)|<\epsilon$ for each $n,m\geq N$ and $x\in [a,b]$ is equivalent to say that $|f_n-f_m|_\infty<\epsilon$ for each $n,m\geq N$.

The condition 2. follows by taking $m\to \infty$, so that $f_m(x)\to f(x)$ and, therefore, we get that
$$
|f_n(x)-f(x)|<\epsilon
$$
for all $n\geq N$ and $x\in [a,b]$.
Finally, this is equivalent to say that
$$
|f_n-f|_{\infty} < \epsilon
$$
for every $n\geq N$ and, crucially, by a fundamental theorem of calculus, this implies that $f$ is continuous because each of the $f_n$ is continuous by assumption.

\section{
$C[a,b]$ is separable
\label{app:B}
}

This is one of the most important examples of a separable Banach space.
The key idea is that every continuous function can be uniformly approximated by polynomials with rational coefficients.
Since there are only countably many such polynomials, they form a countable dense subset of $C[a,b]$.
The proof relies on the Weierstrass approximation theorem.

{\bf Theorem:}
The space
$
C[a,b]
$
equipped with the supremum norm
$$
|f|_\infty
=
\supr{x\in[a,b]}\,\,|f(x)|
$$
is separable.

{\bf Proof:}

{\em Step 1 — A Countable Family of Polynomials.}
Consider the set
$
\mathcal P_{\mathbb Q}
$
of all polynomials with rational coefficients:
$$
p(x)=a_0+a_1x+\cdots+a_nx^n,
\qquad
a_k\in\mathbb Q.
$$
We first show that this set is countable.
For each degree $n$, the set of such polynomials is:
$
\mathbb Q^{n+1}.
$
Since:
\begin{itemize}
\item $\mathbb Q$ is countable,
\item finite products of countable sets are countable,
\end{itemize}
it follows that:
$
\mathbb Q^{n+1}
$
is countable.
Now:
$$
\mathcal P_{\mathbb Q}
=
\bigcup_{n=0}^\infty \mathbb Q^{n+1}.
$$
A countable union of countable sets is countable.
Therefore:
$
\mathcal P_{\mathbb Q}
$
is countable.

{\em Step 2 — Density of Rational Polynomials:}
We now show that:
$$
\overline{\mathcal P_{\mathbb Q}}
=
C[a,b].
$$
Let
$
f\in C[a,b]
$
and let
$
\varepsilon>0.
$
By the Weierstrass approximation theorem, there exists a polynomial
$$
p(x)=a_0+a_1x+\cdots+a_nx^n
$$
with real coefficients such that:
$$
|f-p|_\infty<\frac{\varepsilon}{2}.
$$
Now approximate each real coefficient $a_k$ by a rational number $q_k$ sufficiently accurately.
Define:
$$
q(x)=q_0+q_1x+\cdots+q_nx^n.
$$
Because polynomial evaluation depends continuously on the coefficients, the approximation can be chosen so that:
$$
|p-q|_\infty<\frac{\varepsilon}{2}.
$$
Then:
$$
|f-q|_\infty
\le
|f-p|_\infty+|p-q|_\infty
<
\varepsilon.
$$
Since $q\in\mathcal P_{\mathbb Q}$, we conclude that
$
\mathcal P_{\mathbb Q}
$
is dense in $C[a,b]$.

{\em Conclusion:}
The space $C[a,b]$ contains a countable dense subset.
Therefore:
$
C[a,b]
$
is separable.

{\em Intuition:}
This theorem says something remarkable:
\begin{quotation}
every continuous function can be approximated arbitrarily well using only countably many functions.
\end{quotation}
Even though:
\begin{itemize}
\item $C[a,b]$ is infinite-dimensional,
\item and contains uncountably many functions,
\end{itemize}
its topology is still controlled by a countable set of approximating functions.
This is one of the reasons Fourier analysis and numerical approximation become possible.

\section{
$L^2[a,b]$ is complete
\label{app:C}
}

The completeness of $L^2[a,b]$ is one of the fundamental results of functional analysis.
However, the proof is subtler than for $C[a,b]$, because convergence in $L^2$ is not pointwise convergence and functions in $L^2$ are equivalence classes modulo sets of measure zero.
Still, the core idea is the same:
\begin{quotation}
a Cauchy sequence should converge to an element still inside the space.
\end{quotation}
Let us prove it carefully while keeping the intuition visible.

\subsection{The Space $L^2[a,b]$}

Recall that:
$$
L^2[a,b]
=
\left\{
f:[a,b]\to\mathbb R
:
\int_a^b |f(x)|^2\,dx<\infty
\right\},
$$
where, by definition, {\bf functions differing only on sets of measure zero are identified}.
The norm is:
$$
|f|_2
=
\left(
\int_a^b |f(x)|^2\,dx
\right)^{1/2}.
$$

\subsection{What Must Be Proven?}

We must show:
\begin{quotation}
every Cauchy sequence in $L^2[a,b]$ converges in $L^2$ to another element of $L^2[a,b]$.
\end{quotation}
So let, $\{f_n:n\in \NN\}$ be Cauchy in the $L^2$-norm.
That is
$$
\forall\varepsilon>0,\quad
\exists N
\text{ such that }
m,n\ge N
\Rightarrow
|f_n-f_m|_2<\varepsilon.
$$
Equivalently:
$$
\int_a^b |f_n-f_m|^2\,dx
\to 0.
$$

\subsection{Extract a Rapidly Convergent Subsequence}

Since $\{f_n:n\in \NN\}$ is Cauchy, we can choose a subsequence
$\{f_{n_k}:k\in \NN\}$
such that:
$$
|f_{n_{k+1}}-f_{n_k}|_2
<
2^{-k}.
$$
This is a standard ``rapid convergence'' trick.
Define:
$$
g_k
=
f_{n_{k+1}}-f_{n_k}.
$$
Then:
$$
|g_k|_2<2^{-k}.
$$
Hence:
$$
\sum_{k=1}^\infty |g_k|_2
<
\infty.
$$

\subsection{Build the Candidate Limit}

Consider the series:
$$
\sum_{k=1}^\infty g_k.
$$
Formally:
$$
f_{n_m}
=
f_{n_1}
+
\sum_{k=1}^{m-1} g_k.
$$
The idea is:
\begin{itemize}
\item if the series converges,
\item then the subsequence converges.
\end{itemize}
Now define:
$$
G(x)
=
\sum_{k=1}^\infty |g_k(x)|.
$$
Using standard measure-theoretic arguments (essentially Minkowski/Fatou), one proves:
$$
G(x)<\infty
$$
for almost every $x$.
Thus:
$$
\sum_{k=1}^{\infty} g_k(x)
$$
converges almost everywhere.
Define:
$$
f(x)
=
f_{n_1}(x)
+
\sum_{k=1}^\infty g_k(x).
$$
Then:
$$
f_{n_k}(x)\to f(x)
$$
almost everywhere.

\subsection{Convergence in $L^2$}
Now observe:
$$
f-f_{n_k}
=
\sum_{j=k}^\infty g_j.
$$
Using the triangle inequality:
$$
|f-f_{n_k}|_2
\le
\sum_{j=k}^\infty |g_j|_2.
$$
But:
$$
\sum_{j=k}^\infty 2^{-j}\to 0.
$$
Therefore:
$$
|f-f_{n_k}|_2\to 0.
$$
So the subsequence converges in $L^2$.

\subsection{Finish the Proof}

Since the original sequence $\{f_n:n\in \NN\}$ is Cauchy, and a subsequence converges to $f$, the whole sequence converges to $f$ in $L^2$.
Finally:
$$
|f|_2<\infty,
$$
so:
$$
f\in L^2[a,b].
$$
Thus $L^2[a,b]$ is complete.
Hence it is a Hilbert space.

\subsection{The Key Difference from $C[a,b]$}

For $C[a,b]$:
\begin{itemize}
\item convergence was uniform,
\item limits preserved continuity.
\end{itemize}
For $L^2[a,b]$:
\begin{itemize}
\item convergence is integral-based,
\item functions may oscillate wildly pointwise,
\item only the average square error matters.
\end{itemize}
Thus $L^2$-convergence is much weaker than uniform convergence.

\subsection{Geometric Intuition}

The quantity:
$$
|f-g|_2^2
=
\int_a^b |f-g|^2
$$
measures the ``mean square distance'' between functions.
A Cauchy sequence in $L^2$ means:
\begin{quotation}
The functions become arbitrarily close on average.
\end{quotation}
Completeness guarantees that this approximation process converges to another square-integrable function.
This is exactly why Hilbert spaces are the natural setting for:
\begin{itemize}
\item Fourier series,
\item quantum mechanics,
\item PDEs,
\item signal processing,
\item variational methods.
\end{itemize}

\section{
Separability of $L^2[a,b]$
\label{app:D}
}

The proof is conceptually similar to the separability of $C[a,b]$, but now the approximating objects are not polynomials alone. Instead, one uses:
\begin{enumerate}
\item step functions,
\item rational coefficients,
\item rational interval endpoints.
\end{enumerate}
The key idea is:
\begin{quotation}
Every $L^2$-function can be approximated arbitrarily well by simple functions built from countably many parameters.
\end{quotation}

{\bf Theorem:}
The Hilbert space
$L^2[a,b]$
is separable.

{\bf Proof:}

{\em Step 1 — A Countable Family of Step Functions:}
Consider the set $\mathcal S$ of step functions of the form
$$
s(x)
=
\sum_{k=1}^n q_k\,\chi_{I_k}(x),
$$
where:
\begin{itemize}
\item $q_k\in\mathbb Q$,
\item each $I_k$ is an interval with rational endpoints,
\item $\chi_{I_k}$ denotes the characteristic function of $I_k$.
\end{itemize}
For example:
$$
s(x)
=
2\chi_{[0,1/2)}(x)
-
\frac{3}{5}
\chi_{[1/2,1]}(x).
$$

{\em Step 2 — The Set $\mathcal S$ Is Countable:}
There are only countably many:
\begin{itemize}
\item rational numbers,
\item rational endpoints,
\item finite sequences of such data.
\end{itemize}
Therefore, there are only countably many step functions in $\mathcal S$.
Hence, $\mathcal S$ is countable.

{\em Step 3 — Density of Step Functions:}
We now show that $\mathcal S$ is dense in $L^2[a,b]$.
A fundamental theorem of measure theory states:
\begin{quotation}
Every square-integrable function can be approximated arbitrarily well in $L^2$ by step functions.
\end{quotation}
Thus, for every $f\in L^2[a,b]$ and every $\varepsilon>0$, there exists a step function
$$
s(x)=\sum_{k=1}^n a_k\chi_{J_k}(x)
$$
such that
$$
|f-s|_2<\frac\varepsilon 2.
$$
Here:
\begin{itemize}
\item the coefficients $a_k$ are real,
\item the interval endpoints may be irrational.
\end{itemize}

{\em Step 4 — Approximate the Coefficients and Intervals Rationally:}
Now approximate:
\begin{itemize}
\item each coefficient $a_k$ by a rational number $q_k$,
\item each interval endpoint by rational endpoints.
\end{itemize}
This produces a rational step function $r\in\mathcal S$.
Because the $L^2$-norm depends continuously on:
\begin{itemize}
\item the coefficients,
\item and the interval lengths,
\end{itemize}
the approximations can be chosen sufficiently accurately so that
$$
|s-r|_2<\frac\varepsilon2.
$$
Then, by the triangle inequality,
$$
|f-r|_2
\le
|f-s|_2+|s-r|_2
<
\varepsilon.
$$
Thus:
$$
\mathcal S
$$
is dense in $L^2[a,b]$.

{\em Conclusion:}
We have constructed a:
\begin{itemize}
\item countable,
\item dense
\end{itemize}
subset of $L^2[a,b]$.
Therefore, $L^2[a,b]$ is separable.

{\em Intuition:}
This theorem says:
\begin{quotation}
Even though $L^2[a,b]$ contains enormously complicated objects, every element can still be approximated using only countably many simple building blocks.
\end{quotation}
This is one of the reasons:
\begin{itemize}
\item Fourier series,
\item numerical approximations,
\item spectral decompositions,
\item quantum mechanics,
\end{itemize}
become mathematically manageable.

{\em Connection with Hilbert Spaces:}
Separability is particularly important for Hilbert spaces because it allows the existence of countable orthonormal bases:
$$
f
=
\sum_{n=1}^\infty c_n e_n.
$$
Without separability, one would generally need uncountably many basis vectors, and generalized Fourier analysis would become much more complicated.
This is one of the reasons most Hilbert spaces appearing in physics are separable.

%
%\section{
%Compactness in infinite dimension
%\label{app:A}
%}
%
%% https://chatgpt.com/share/6a0c5541-7498-83e9-a94f-8ca91aa5e791
%
%The compactness of a subset $K$ of a metric space $X$ only makes sense if $X$ is finite-dimensional?
%No. Compactness is a fundamental topological notion and makes perfect sense in {\em any} metric space, including infinite-dimensional ones.
%For a subset $K \subset X$ of a metric space $(X,d)$, compactness means: 
%\begin{quote}
%Every open cover of $K$ admits a finite subcover.
%\end{quote}
%In metric spaces, this is equivalent to: 
%\begin{quote}
%Every sequence in $K$ has a convergent subsequence whose limit lies in $K$.
%\end{quote}
%This definition does not depend on dimension at all.
%What {\em does} change drastically between finite- and infinite-dimensional spaces is {\bf which sets are compact}.
%
%\subsection{Finite-dimensional case}
%
%In $\mathbb{R}^n$, the Heine–Borel theorem says:
%$$
%K \subset \mathbb{R}^n
%\text{ is compact }
%\iff
%K \text{ is closed and bounded.}
%$$
%This is a very special property of finite-dimensional normed spaces.
%
%\subsection{Infinite-dimensional case}
%
%In infinite-dimensional normed spaces, closed and bounded sets are {\em usually not compact}.
%For example, consider the Hilbert space
%$$
%\ell^2
%=
%\left\{
%x=(x_1,x_2,\dots):
%\sum_{n=1}^\infty |x_n|^2 < \infty
%\right\}
%.
%$$
%Its closed unit ball
%$$
%B=\{x\in \ell^2:|x|\le 1\}
%$$
%is closed and bounded, but {\bf not compact}.
%A standard proof uses the canonical basis vectors
%$$
%e_1=(1,0,0,\dots),\quad
%e_2=(0,1,0,\dots),\quad \dots
%$$
%which satisfy
%$$
%|e_n-e_m|=\sqrt{2}
%\quad (n\neq m).
%$$
%Hence no subsequence can be Cauchy, so no subsequence converges. Therefore $B$ is not compact.
%
%\subsection{What replaces Heine–Borel?}
%
%In metric spaces, compactness is closely related to {\bf total boundedness}.
%
%For metric spaces:
%
%$$
%K \text{ compact}
%\iff
%K \text{ complete and totally bounded}.
%$$
%
%Total boundedness is stronger than ordinary boundedness.
%A set is totally bounded if:
%\begin{quote}
%For every $\varepsilon>0$, it can be covered by finitely many balls of radius $\varepsilon$.
%\end{quote}
%In finite dimensions, boundedness already implies total boundedness. In infinite dimensions, it generally does not.
%
%\subsection{Intuition}
%
%Infinite-dimensional balls contain ``too many independent directions.''
%The sequence $e_n$ in $\ell^2$ keeps jumping into new orthogonal directions and never accumulates anywhere.
%That phenomenon cannot happen inside bounded subsets of finite-dimensional spaces.
%So compactness absolutely makes sense in infinite-dimensional spaces — it simply becomes much more restrictive.


\end{document}
%
% ****** End of file apstemplate.tex ******


